% A4縦・1カラム．参考PDFに近い余白，章見出し，右上ページ番号．
\usepackage[haranoaji]{luatexja-preset}
\usepackage[top=30mm,bottom=30mm,left=26mm,right=26mm,headheight=16pt,headsep=10mm]{geometry}
\usepackage{graphicx}
\usepackage{amsmath,amssymb}
\usepackage{booktabs,tabularx}
\usepackage{xcolor}
\usepackage{url}
\usepackage{fancyhdr}
\usepackage[unicode,hidelinks,bookmarksnumbered]{hyperref}

\setcounter{secnumdepth}{2}
\setcounter{tocdepth}{2}
\renewcommand{\listfigurename}{図目次}
\renewcommand{\listtablename}{表目次}
\renewcommand{\bibname}{参考文献}
\renewcommand{\baselinestretch}{1.12}
\setlength{\parindent}{1\zw}
\setlength{\parskip}{0pt}
\raggedbottom

\fancypagestyle{thesis}{%
  \fancyhf{}
  \fancyhead[R]{\small\thepage}
  \renewcommand{\headrulewidth}{0pt}
  \renewcommand{\footrulewidth}{0pt}
}
% 章の開始ページにも同じ位置に番号を置く．
\fancypagestyle{plain}{%
  \fancyhf{}
  \fancyhead[R]{\small\thepage}
  \renewcommand{\headrulewidth}{0pt}
  \renewcommand{\footrulewidth}{0pt}
}
\makeatletter
\let\ps@plainhead\ps@plain
\let\ps@plainfoot\ps@plain
\makeatother

% 執筆時に，この命令全体を実際の本文へ置き換える．
% 記入案内は完成した研究内容ではないため，灰色で表示する．
\newcommand{\writehere}[1]{\par\noindent{\color{gray}\small［記入内容：#1］}\par}
\newcommand{\unnumberedchapter}[1]{%
  \chapter*{#1}%
  \phantomsection
  \addcontentsline{toc}{chapter}{#1}%
  \markboth{#1}{#1}%
}
